\documentclass[border=0pt]{standalone}
% The shared preamble of every docs diagram (docs/diagrams/<name>.tex).
% dvisvgm's driver first, so patterns and opacity survive into the SVG.
\def\pgfsysdriver{pgfsys-dvisvgm.def}
\usepackage[siunitx,RPvoltages]{circuitikz}
\usepackage{xcolor}
% Latin Modern: outline fonts for every glyph a label uses, µ included.
\usepackage[T1]{fontenc}
\usepackage{lmodern}
\usetikzlibrary{arrows.meta,backgrounds,calc,patterns,positioning}
\tikzset{
  % A white card, so diagrams read in mdBook's dark themes too.
  every picture/.append style={show background rectangle,
    background rectangle/.style={fill=white,inner sep=4pt}},
  % Layout geometry: layers told apart by fill and hatching, never colour alone.
  layer a/.style={draw=black,line width=0.6pt,fill=gray!20},
  layer b/.style={draw=black,line width=0.6pt,fill=gray!45},
  layer c/.style={draw=black,line width=0.6pt,fill=white,postaction={pattern=north east lines}},
  layer d/.style={draw=black,line width=0.6pt,fill=white,postaction={pattern=crosshatch dots}},
  marker/.style={draw=black,dashed,line width=0.6pt},
  % A measured distance and its label.
  dim/.style={{Latex[length=2mm]}-{Latex[length=2mm]},line width=0.5pt},
  dim label/.style={font=\small,fill=white,inner sep=1pt},
  % Where a rule fails: thick, dashed and labelled, not only red.
  violation/.style={draw=red!75!black,line width=1.4pt,dash pattern=on 3pt off 2pt},
  violation label/.style={font=\small\bfseries,text=red!75!black},
  caption/.style={font=\small},
}

\begin{document}
\begin{circuitikz}[x=0.3cm,y=0.45cm]
  % Cross-section of the example (µm): nwell 0-10 with ntap 1-2 on VDD and
  % pdiff 4-7 on VDDIO; ndiff victim 12-20; ptap 30-31.
  \fill[layer a] (-1,-11) rectangle (33,0);
  \fill[layer b] (0,-6) rectangle (10,0);
  \fill[layer d] (1,-1) rectangle (2,0);
  \fill[layer c] (4,-1) rectangle (7,0);
  \fill[layer d] (12,-1) rectangle (20,0);
  \fill[layer c] (30,-1) rectangle (31,0);
  \draw (1.5,0) -- (1.5,1.2) node[caption,anchor=south east] {VDD \SI{1.8}{\volt}};
  \draw (5.5,0) -- (5.5,1.2) node[caption,anchor=south west] {VDDIO \SI{3.3}{\volt}};
  \node[caption,anchor=north] at (1.5,-1.1) {ntap};
  \node[caption,anchor=south] at (16,0.2) {ndiff (victim)};
  \node[caption,anchor=south] at (30.5,0.2) {ptap};
  \node[caption,anchor=west] at (0.2,-5.3) {nwell};
  \node[caption,anchor=west] at (22,-10.3) {p substrate};
  % The thyristor: PNP (pdiff, nwell, substrate), NPN (ndiff, substrate, nwell).
  \node[pnp,anchor=E] (P) at (5.5,-1) {};
  \node[npn,anchor=E,yscale=-1] (N) at (16,-1) {};
  \draw (P.C) -- ++(0,-4.5) coordinate (pc) -| ($(N.B)+(-2,0)$) -- (N.B);
  \draw (N.C) -- ++(0,-6) -| (2.5,0 |- P.B) -- (P.B);
  \node[caption,anchor=west] at ($(P.C)+(0.5,0.4)$) {PNP};
  \node[caption,anchor=west] at ($(N.C)+(0.5,-0.4)$) {NPN};
  % The victim's distance to its tap.
  \draw[dim] (12,3.4) -- node[dim label,above] {\SI{18}{\micro\metre} to ptap, limit \SI{5}{\micro\metre}} (30,3.4);
  \draw[densely dotted] (12,0) -- (12,3.4) (30,0) -- (30,3.4);
  \draw[violation] (12,0) circle[radius=5pt];
  \node[violation label,anchor=west,align=left] at (22,-3)
    {pdiff can rise \SI{1.5}{\volt}\\ above its well:\\ an aggressor};
\end{circuitikz}
\end{document}
